% ============================================================================
%  光储充微电网物联网平台 — 技术方案（Python 轻量版）
%  版本: V1.0 | 日期: 2026年6月
% ============================================================================
\documentclass[12pt,a4paper]{ctexart}
\usepackage{xeCJK}
\usepackage{graphicx,float,fancyhdr,hyperref,geometry,longtable,booktabs,enumitem,caption,xcolor,tikz}
\setmainfont{Noto Serif SC}
\setCJKmainfont{SimSun}
\setCJKsansfont{SimHei}
\geometry{a4paper,margin=2.5cm,top=3cm,bottom=2.5cm}
\pagestyle{fancy}
\fancyhead[L]{光储充微电网物联网平台·技术方案}
\fancyhead[R]{V1.0}
\fancyfoot[C]{\thepage}
\hypersetup{colorlinks=true,linkcolor=blue,pdftitle={光储充微电网物联网平台技术方案}}

\begin{document}

% 封面
\begin{titlepage}
\centering\vspace*{3.5cm}
{\Huge\bfseries\sffamily 光储充微电网物联网平台}\\[0.8cm]
{\LARGE\bfseries 技术方案}\\[1.5cm]
{\large 轻量版 · Python 全开源}\\[2.5cm]
{\large\begin{tabular}{rl}
\textbf{版 本 号:} & V1.0 \\
\textbf{文档状态:} & 初稿 \\
\textbf{日　　期:} & 2026年6月 \\
\end{tabular}}\\[1.5cm]
{\normalsize 迪格(杭州)物联科技有限公司}
\end{titlepage}

\newpage\tableofcontents\newpage

% ===== 一、项目概述 =====
\section{项目概述}

\subsection{项目背景}
客户需要一套\textbf{可源码交付}的物联网平台，用于光储充微电网场站的数据采集与监控。客户拿到源码后自行掌握二次开发、部署和运维。

\textbf{我方自有物联网平台与本项目独立}——本项目交付的是基于 Python 全开源技术栈的轻量级物联网平台，专为满足客户源码交付需求而构建。

\subsection{建设目标}
\begin{enumerate}[leftmargin=*]
  \item \textbf{四协议统一接入}：Modbus RTU、Modbus TCP、IEC 104、OPC UA 实际跑通
  \item \textbf{时序+关系双存储}：TDengine（时序）+ PostgreSQL（关系），开源免费
  \item \textbf{2D组态监控}：HTML5 Canvas 实现，浏览器直接访问，无需插件
  \item \textbf{数据推送}：MQTT / HTTP 实时推送至第三方应用
  \item \textbf{源码交付}：完整 Python 源码，客户可自行编译部署和二次开发
\end{enumerate}

\subsection{项目范围与分工}
\textbf{我方负责物联网平台层——平台部署、协议适配、组态配置、数据推送，确保系统实际跑通。}

\begin{table}[H]\centering
\begin{tabular}{p{3.5cm}p{5cm}p{5cm}}
\toprule\textbf{工作内容} & \textbf{我方（平台层）} & \textbf{客户（应用层）} \\
\midrule
数据采集引擎 & \checkmark\ 4协议采集，实际跑通 & — \\
2D组态 & \checkmark\ 组态页面+数据绑定 & \checkmark\ 业务画面定制 \\
数据库 & \checkmark\ TDengine+PG 部署+表结构 & — \\
数据推送 & \checkmark\ MQTT/HTTP 通道开发 & \checkmark\ 第三方对接 \\
碳排优化 & \checkmark\ 数据采集+核算模型 & \checkmark\ 策略制定 \\
源码处置 & — & \checkmark\ 客户自行处置 \\
\bottomrule
\end{tabular}
\end{table}

% ===== 二、技术选型 =====
\section{技术选型}

\subsection{选型策略}

本项目定位为\textbf{轻量学习版}，预算约 5 万元，技术选型原则：
\begin{itemize}
  \item \textbf{全开源}：所有组件开源免费，无 License 成本
  \item \textbf{Python 优先}：生态丰富、开发高效、客户易于接手
  \item \textbf{异步架构}：asyncio + FastAPI，支撑多设备并发采集
  \item \textbf{降级模式}：无 PostgreSQL/TDengine 时自动降级为 SQLite，单机即跑
\end{itemize}

\subsection{技术栈}

\begin{table}[H]\centering
\begin{tabular}{p{3cm}p{4.5cm}p{5cm}}
\toprule\textbf{层次} & \textbf{选型} & \textbf{说明} \\
\midrule
核心语言 & Python 3.10+ & 异步非阻塞，生态丰富 \\
Web框架 & FastAPI + uvicorn & 自动 API 文档，高性能异步 \\
Modbus & pymodbus 3.6+ & RTU + TCP 双模式 \\
IEC 104 & 自研轻量客户端 & ASDU 解析，总召+变化传输 \\
OPC UA & opcua-asyncio & 订阅+轮询双模式 \\
时序存储 & TDengine 3.x & 列式存储，10-20:1 压缩 \\
关系存储 & PostgreSQL 15+ & SQLAlchemy async ORM \\
消息推送 & paho-mqtt + httpx & MQTT+HTTP 双通道 \\
前端 & HTML5 Canvas + Vanilla JS & 零依赖 2D 组态 \\
部署 & Docker Compose & 一键部署全部组件 \\
\bottomrule
\end{tabular}
\end{table}

% ===== 三、总体架构 =====
\section{总体架构设计}

系统采用四层架构：

\begin{enumerate}[leftmargin=*]
  \item \textbf{设备接入层}：协议适配器（Modbus RTU/TCP、IEC 104、OPC UA），每个设备一个适配器实例，独立异步采集循环
  \item \textbf{核心服务层}：采集调度引擎（CollectorEngine）+ 告警引擎（AlarmEngine）+ 推送引擎（PushEngine），通过回调链串联
  \item \textbf{数据存储层}：TDengine 存遥测时序数据，PostgreSQL 存设备档案/配置/告警，降级模式自动切 SQLite
  \item \textbf{Web 服务层}：FastAPI REST API + WebSocket 实时推送 + 2D 组态静态页面
\end{enumerate}

\subsection{数据流向}

采集流程：\textbf{调度器定时触发 → 协议适配器读取硬件 → 解析为 PointValue → 并行写入 TDengine → 转发告警引擎和推送引擎 → WebSocket 广播至前端}

% ===== 四、协议方案 =====
\section{数据采集协议方案}

\subsection{Modbus RTU/TCP}

基于 pymodbus 异步客户端实现。RTU 模式通过 RS-485 串口轮询，TCP 模式通过以太网并发连接。支持功能码 01-04、06、16，自动解析 int16/uint16/int32/float32/float64 数据类型。每个设备独立采集周期（1-60秒可配），采集失败自动重连。

\subsection{IEC 60870-5-104}

轻量自研 IEC 104 客户端（主站模式），实现：
\begin{itemize}
  \item APCI 帧解析（I 帧发送序列号管理、U 帧 STARTDT/STOPDT）
  \item ASDU 类型支持：M\_ME\_NA\_1(9)、M\_ME\_NB\_1(11)、M\_ME\_NC\_1(13)、M\_SP\_NA\_1(1)、M\_DP\_NA\_1(3)
  \item 总召命令（C\_IC\_NA\_1）+ 时钟同步（C\_CS\_NA\_1）
  \item TCP/2404 连接，自动重连
\end{itemize}

\subsection{OPC UA}

基于 opcua-asyncio，支持：
\begin{itemize}
  \item Browse 地址空间自动发现变量节点
  \item 订阅（Subscription）模式：数据变化主动推送
  \item 轮询（Read）模式：周期读取
  \item 安全策略 None / Basic256Sha256 + 用户名密码认证
\end{itemize}

% ===== 五、应用功能 =====
\section{应用功能设计}

\subsection{设备与点位管理}
\begin{itemize}
  \item 设备档案：设备名称、类型（光伏/储能/充电桩/电表/传感器）、协议、通讯参数
  \item 点位配置：协议地址、数据类型、采集周期、缩放系数、告警阈值
  \item 批量导入：支持 Excel/CSV 点位模板批量导入
\end{itemize}

\subsection{2D组态监控}

基于 HTML5 Canvas 的 2D 组态画面：
\begin{itemize}
  \item \textbf{画面}：场站总览、光伏监控、储能监控、充电桩监控、碳排仪表盘、告警总览（6 画面）
  \item \textbf{图元}：光伏阵列、逆变器、PCS、变压器、充电桩、电网接口等电力设备图标
  \item \textbf{实时交互}：点击设备弹出数据详情面板，告警变色闪烁，WebSocket 实时刷新
\end{itemize}

\subsection{碳排优化细则}
\label{sec:carbon}

\textbf{碳排放源识别}：光储充场站无直接排放源，间接排放来自电网购电（范围二）。碳减排来自光伏替代电网电量和充电桩替代燃油。

\textbf{数据采集要求}：从逆变器采集光伏发电量（15分钟）、PCS采集充放电量（15分钟）、关口表采集电网购售电量（15分钟）、充电桩采集服务电量（累计）、辐照仪采集辐照度（1分钟）。

\textbf{核算标准}：采用 GB/T 32150 方法学，电网排放因子引用生态环境部最新发布值。\\
净碳排放 = 电网购电量 × 排放因子 − 光伏发电量 × 排放因子 − 充电量 × 燃油替代因子

\textbf{报告周期}：实时仪表盘（1分钟刷新）+ 日报 + 月报 + 年报

\subsection{关键KPI}

\begin{table}[H]\centering
\begin{tabular}{p{2cm}p{3.5cm}p{2.5cm}p{2cm}}
\toprule\textbf{类别} & \textbf{指标} & \textbf{目标值} & \textbf{周期} \\
\midrule
采集质量 & 数据完整率 & $\geq$ 99\% & 日/月 \\
采集质量 & 设备在线率 & $\geq$ 99\% & 日/月 \\
采集质量 & 平均采集延迟 & $\leq$ 3秒 & 实时 \\
平台性能 & 数据推送延迟 & $\leq$ 1秒 & 实时 \\
平台性能 & API响应时间 & $\leq$ 500ms & 实时 \\
碳排指标 & 光伏碳减排量 & 按装机核算 & 月/年 \\
碳排指标 & 场站净碳排放 & $\leq$ 0(碳中和) & 月/年 \\
碳排指标 & 可再生能源占比 & $\geq$ 80\% & 月 \\
\bottomrule
\end{tabular}
\end{table}

\subsection{告警管理}

阈值告警（高限/上上限/低限/下下限）、变化率告警、设备离线告警。三级分类 P0/P1/P2。WebSocket 实时推送至前端 + MQTT 推送至第三方。

\subsection{数据推送}

MQTT（默认）+ HTTP POST 双通道。消息格式 JSON，规则引擎按设备类型/点位过滤。断网本地缓存、恢复补推。

% ===== 六、存储方案 =====
\section{数据存储方案}

\subsection{TDengine 时序存储}
超级表 device\_telemetry（ts + value + quality），TAGs（device\_id, point\_id, point\_name, unit, device\_type）。每个设备点位独立子表。列式存储+差值编码，压缩比 10-20:1。1分钟级数据保留 12 个月。

\subsection{PostgreSQL 关系存储}
存储设备档案（devices）、点位配置（data\_points）、告警记录（alarm\_records）、推送目标（push\_targets）。SQLAlchemy async ORM + Alembic 迁移。

\subsection{降级模式}
无 TDengine/PostgreSQL 环境时，自动降级为 SQLite 存储全部数据，单机即可完整运行——适合开发测试和轻量部署场景。

% ===== 七、部署方案 =====
\section{部署方案}

\subsection{服务器配置（$\leq$500设备）}

\begin{table}[H]\centering
\begin{tabular}{p{3cm}p{3cm}p{4cm}}
\toprule\textbf{资源} & \textbf{最低} & \textbf{推荐} \\
\midrule
CPU & 4 核 & 8 核 \\
内存 & 8 GB & 16 GB \\
磁盘 & 200 GB SSD & 500 GB SSD \\
系统 & Ubuntu 20.04+ & Ubuntu 22.04 LTS \\
\bottomrule
\end{tabular}
\end{table}

\subsection{Docker 一键部署}

\begin{verbatim}
cd scripts && docker-compose up -d
\end{verbatim}

一键启动 IoT 平台 + PostgreSQL + TDengine + EMQ X 四个容器。

\subsection{裸机部署}

\begin{verbatim}
pip install -r requirements.txt
python scripts/init_db.py
python run.py
# 访问 http://localhost:8000/docs
\end{verbatim}

\subsection{数据库说明}
数据库采用\textbf{开源免费版本}：TDengine Community（AGPL v3）+ PostgreSQL（PostgreSQL License）。客户按《部署手册》自行安装，不产生 License 费用。

% ===== 八、实施计划 =====
\section{实施计划}

\begin{table}[H]\centering
\begin{tabular}{p{1.2cm}p{3cm}p{1.5cm}p{5cm}}
\toprule\textbf{阶段} & \textbf{工作内容} & \textbf{工期} & \textbf{验收标准} \\
\midrule
P1 & 需求确认+环境准备 & 1周 & 设备清单、通讯参数、服务器就绪 \\
P2 & 平台部署+数据库初始化 & 1周 & 平台可登录，数据库连接正常 \\
P3 & 四协议适配+联调 & 3周 & \textbf{各协议至少接入1台真实设备跑通} \\
P4 & 2D组态+数据推送 & 2周 & 组态画面实时数据展示，推送链路正常 \\
P5 & 联调测试+文档 & 2周 & 端到端测试通过，数据准确率达标 \\
P6 & 验收+培训 & 1周 & 交付物齐全，客户签字验收 \\
\hline
\textbf{合计} & & \textbf{10周} & \\
\bottomrule
\end{tabular}
\end{table}

\subsection{交付物}
\begin{itemize}
  \item IoT 平台完整 Python 源码（含 Git 仓库）
  \item 技术方案（本文档）+ 部署手册 + 协议接入指南 + API 文档
  \item 数据库初始化脚本 + Docker Compose 部署文件
  \item 验收测试报告
\end{itemize}

% ===== 九、运维服务 =====
\section{运维服务方案}

\textbf{服务范围}：保障交付的 IoT 平台长期稳定运行——四协议采集链路、TDengine/PostgreSQL 数据库运行、数据推送链路、2D 组态可用性。

\begin{table}[H]\centering
\begin{tabular}{p{2cm}p{2cm}p{4cm}p{4cm}}
\toprule\textbf{级别} & \textbf{响应} & \textbf{场景} & \textbf{方式} \\
\midrule
P0 紧急 & $\leq$30min & 平台宕机、全站数据中断 & 远程+电话 \\
P1 重要 & $\leq$2h & 单协议中断、推送断开 & 远程 \\
P2 一般 & $\leq$8h & 个别设备离线 & 工单 \\
\bottomrule
\end{tabular}
\end{table}

\textbf{不含}：平台版本升级（交付版本长期可用，升级迭代不纳入本项目）；客户应用层运维；服务器硬件。

% ===== 十、投资概算 =====
\section{投资概算（脱敏版）}

\subsection{费用构成}

\begin{table}[H]\centering
\begin{tabular}{p{1.2cm}p{3.5cm}p{3cm}p{5cm}}
\toprule\textbf{序号} & \textbf{费用项目} & \textbf{计价} & \textbf{说明} \\
\midrule
1 & 四协议适配开发 & 按工作量 & Modbus RTU/TCP + IEC104 + OPCUA 采集引擎 \\
2 & 2D组态开发 & 6个画面 & 总览/光伏/储能/充电桩/碳排/告警 \\
3 & 数据推送开发 & 3通道 & MQTT + HTTP + WebSocket \\
4 & 数据库部署 & 2个 & TDengine + PostgreSQL 建库建表 \\
5 & 平台部署+联调 & 按人天 & 环境部署+协议调试+性能调优 \\
6 & 文档+培训 & 按项 & 部署手册+API文档+培训 \\
\hline
\multicolumn{3}{r}{\textbf{合计}} & \textbf{详见商务报价函} \\
\bottomrule
\end{tabular}
\end{table}

\textbf{说明}：
\begin{itemize}
  \item 全开源组件，无第三方 License 费用
  \item 数据库（TDengine Community + PostgreSQL）免费，客户自装
  \item 不含服务器硬件，客户自行采购
  \item 年度维保服务按年另签
\end{itemize}

\subsection{付款方式建议}
合同签订后→平台部署完成→联调测试通过→终验合格，分四期支付。

\end{document}
